Coherent elastic neutrino-nucleus scattering (CEvNS) offers a compact, Standard-Model-anchored probe of reactor antineutrinos, but the nuclear recoils it produces in germanium at reactor energies are only tens to hundreds of eV, below the reach of conventional ionization or phonon detectors. First predicted by Freedman~\cite{freedman1974} and proposed as a neutrino-detection channel by Drukier and Stodolsky~\cite{drukier1984}, CEvNS was observed only in 2017 by the COHERENT experiment at a stopped-pion source~\cite{akimov2017}, which fixed the cross-section normalization against the electroweak Standard Model. Reactor sources move the target recoil energies an order of magnitude lower, into the tens-of-eV regime, in exchange for an enormous $\bar\nu_e$ flux. A dedicated effort has followed---CONUS+~\cite{conus2025}, Dresden-II~\cite{dresden2022}, and NUCLEUS~\cite{nucleus}---each pushing the deposited-energy threshold toward and below $100$~eV, where the germanium reactor-CEvNS rate is benchmarked by the calculation of Billard~\emph{et~al.}~\cite{billard2017}. The central experimental obstacle is not statistics but threshold: recovering the reactor-CEvNS spectrum demands sensitivity to nuclear recoils that terminate at a few tens of eV.

Quantum Parity Detectors (QPDs) are a superconducting readout concept aimed squarely at this regime~\cite{ramanathan2026}. A QPD infers a deposited energy from the population of quasiparticles it liberates in a superconducting sensor, sensed through the parity-dependent tunneling of individual quasiparticles across a Josephson junction. Because each tunneling event is counted, the scheme offers a path to sub-keV and even eV-scale thresholds on a fieldless phonon calorimeter. It carries, however, an intrinsic limitation: the tunneling signal is bandwidth-limited, and the readout response \emph{saturates} once the instantaneous tunneling rate exceeds the resolving ceiling---here $\sim$$25$~kHz, set by a $\tau_d\simeq40$~$\mu$s dead time. A large deposit drives the burst rate far above this ceiling, so the number of \emph{recorded} tunneling events, and hence the inferred energy, no longer tracks the energy actually deposited. The decisive observable for such a detector is therefore not the deposited energy but the \emph{reconstructed} energy $E_{\rm rec}$ that survives the bandwidth-limited counting chain. A background that is innocuous in deposited energy can, after reconstruction, land squarely on top of an eV-scale signal.

This paper computes that reconstructed-energy observable end to end. We build a forward model that maps the deposited-energy spectrum of the reactor-CEvNS signal and of its surface backgrounds into reconstructed-energy differential rate spectra $dR/dE_{\rm rec}$ for a thin single-sided germanium wafer read out by QPDs. The backgrounds are of three kinds: external electron-recoil continua from through-going cosmic-ray muons and the radiogenic environmental-gamma field; the external nuclear-recoil field of ambient fast neutrons, which scatter elastically off germanium and mimic the CEvNS signature directly; and the internal Ge-intrinsic decays of the cosmogenically activated isotopes $^{3}$H, $^{68}$Ge, and $^{65}$Zn, which deposit from inside the crystal and which no external shield can reach. We treat two quasiparticle-trapping designs, Ta$\to$Al and Al$\to$Hf~\cite{ramanathan2026}, and we place every channel on a single unified phonon energy scale with no ionization quenching, as befits a fieldless calorimeter, so that signal and backgrounds are directly co-added and folded through a common response. Every rate is anchored to a published value or evaluated data library rather than assumed: the deposited CEvNS rate reproduces the Billard germanium benchmark to within $2.4\%$, the cosmic-muon flux matches the sea-level value tabulated by the Particle Data Group to within its $\sim$$30\%$ inter-experiment spread~\cite{pdg}, the environmental-gamma Compton edges are fixed exactly by scattering kinematics, the ambient neutron flux is anchored to the Gordon sea-level measurement~\cite{gordon2004}, and the cosmogenic production rates to the SuperCDMS and EDELWEISS germanium-activation measurements~\cite{amman2019,edelweiss2017}. Because the external channels can be attenuated but the internal floor cannot, we present the background budget in an assumed-shield scenario---external muon, gamma, and neutron fields suppressed at the performance level demonstrated by NUCLEUS~\cite{nucleus}, the Ge-intrinsic floor carried unsuppressed. The response chain that carries these deposited spectra into $E_{\rm rec}$---the sensor-sharing partition, the tunneling-burst statistics, and the non-paralyzable censoring at the $25$~kHz ceiling---is the methodological core of the work.

The headline finding is a qualitative separation on the reconstructed axis. Because the reconstructed energy is calibrated to \emph{estimate} the deposited energy, the reactor-CEvNS signal, whose recoils lie entirely below the saturation onset, reconstructs to tens of eV, with a median near $55$~eV and about two-thirds of its counts below $100$~eV. It reconstructs onto the very lowest end of the observable axis---but it does not sit there alone. The ambient neutron-elastic recoils occupy the same tens-of-eV region, and even after the assumed shield they are the largest term in the reconstructed $10$--$100$~eV region of interest, at $\sim$$4.8\times10^{3}$ counts~kg$^{-1}$day$^{-1}$ against a CEvNS signal of $\sim$$63$: the flagship channel is signal-limited by a nuclear-recoil background that shares its kinematics, not by the readout. The cosmic-muon channel, by contrast, deposits $\sim$$1.5$--$197$~MeV and drives the readout deep into saturation, so its entire deposited spectrum collapses onto a single tens-of-keV plateau ($\approx$$37.6$~keV for Ta$\to$Al, $\approx$$29.9$~keV for Al$\to$Hf), well clear of the signal band. Beneath everything lies the Ge-intrinsic cosmogenic floor, whose $^{68}$Ge and $^{65}$Zn electron-capture lines reconstruct into the shoulder just above the RoI and whose $^{3}$H continuum underlies it. We stress at the outset what this separation is and is not. It is a stage-1 forward-model, or feasibility, calculation, not a sensitivity or discovery claim, and we quote no detection significance. The muon pile-up is an \emph{instrumental artifact} of the modeled bandwidth ceiling, not a physical spectral line, and the shape of the saturated-regime response has no published anchor; the neutron channel carries only an order-of-magnitude accuracy label and its assumed suppression is a modeling input, not a measured rejection; the absolute environmental-gamma normalization is a site-dependent estimate good to a factor of $\sim$$2$; and the crossover between the linear and saturated regimes rides on two low-confidence sensor-sharing parameters, so we report it as a band spanning roughly $50$~eV to $13$~keV rather than a single number. These modeling uncertainties are quantified channel by channel throughout, so that the separation stands on stated assumptions rather than hidden ones.

The remainder of the paper follows the pipeline in order. Section~\ref{sec:model} fixes the detector geometry, the unified phonon energy scale, and the per-sensor tunneling chain. Section~\ref{sec:cevns} builds the deposited reactor-CEvNS signal from an absolute antineutrino flux and the coherent cross section, and anchors it to the Billard benchmark. Section~\ref{sec:backgrounds} computes the deposited cosmic-muon, environmental-gamma, ambient-neutron, and Ge-intrinsic cosmogenic spectra. Section~\ref{sec:response} constructs the bandwidth-limited response matrix $R(E_{\rm rec}\,|\,E_{\rm dep})$ that is the heart of the forward model. Section~\ref{sec:results} folds every channel through it to obtain the reconstructed spectra of Fig.~\ref{fig:spectra} in the assumed-shield scenario, and Secs.~\ref{sec:discussion} and~\ref{sec:conclusions} discuss the implications and limits of the result.
